\documentclass[12pt,letterpaper]{article}
\usepackage[margin=1in]{geometry}
\usepackage[T1]{fontenc}
\usepackage{amsmath,amssymb}
\usepackage{newtxtext,newtxmath}
\usepackage{helvet}
\usepackage{enumitem}
\usepackage{caption,booktabs}
\usepackage[hidelinks]{hyperref}
\captionsetup{font=small,labelfont=bf,skip=4pt}
\renewcommand{\arraystretch}{1.15}
\setlength{\emergencystretch}{2em}
\pagestyle{empty}
\setlength{\parindent}{0pt}
\setlength{\parskip}{0pt}
\AtBeginDocument{\fontsize{12}{13.8}\selectfont}
\setlist[enumerate,1]{label=\arabic*.,leftmargin=0.25in,
labelsep=0.08in,align=left,topsep=0pt,itemsep=14pt,parsep=6pt}

\hypersetup{pdfauthor={Anonymous contributor},pdftitle={ECEN 474/704 - Lab 5 Prelab},pdfsubject={Anonymized laboratory prelab}}
\begin{document}
{\fontsize{11}{12.65}\selectfont
\hfill September 29, 2026\par}
{\centering
{\fontsize{14}{16.1}\selectfont\sffamily\bfseries ECEN 474/704\par}
{\fontsize{14}{16.1}\selectfont Lab 5: Prelab\par}
{\fontsize{14}{16.1}\selectfont Anonymous contributor\par}
}
\vspace{12pt}
\begin{enumerate}

\item \textbf{Comparison of current-mirror topologies.}
Make a table listing the three current-mirror topologies described in
this lab. Rate each topology using good, medium, and bad for output
resistance, accuracy, complexity, and compliance voltage.

\textbf{Solution.}

The comparison assumes comparable device parameters and operating
currents. ``Good'' denotes a desirable
property: high output resistance, high accuracy, low complexity, or
low required compliance voltage.

\begin{center}
\small
\setlength{\tabcolsep}{5pt}
\begin{tabular}{@{}lcccc@{}}
\toprule
Topology & $R_{\mathrm{out}}$ & Accuracy & Complexity & Compliance voltage\\
\midrule
Simple & Bad & Medium & Good & Good\\
Cascode & Good & Good & Medium & Bad\\
Low-voltage cascode & Good & Good & Bad & Medium\\
\bottomrule
\end{tabular}
\end{center}

\textbf{(a) Simple current mirror.}
The two-transistor mirror has the lowest complexity and requires
approximately one overdrive voltage at its output. Its output
resistance is approximately $r_o$. Differences between the reference
and output drain voltages produce current error through
channel-length modulation.

\textbf{(b) Cascode current mirror.}
The cascode devices reduce changes in the drain voltages of the
mirroring devices, improving current accuracy and increasing output
resistance to approximately $g_mr_o^2$ when $g_mr_o\gg1$.
For equal overdrives, the conventional topology requires approximately
$V_{TH}+2V_{OV}$ of output voltage.

\textbf{(c) Low-voltage cascode current mirror.}
A separate bias branch sets the cascode gate voltage, allowing high
output resistance and good current matching with an ideal minimum
compliance voltage near $2V_{OV}$.
This is lower than the conventional cascode requirement but higher
than that of the simple mirror. The additional bias transistor and
bias current increase circuit complexity. Saturation margin increases
the practical compliance voltage above the ideal minimum.

\newpage
\item \textbf{Simple 1:1 current-mirror design.}
Design a simple 1:1 current mirror with a compliance voltage between
$125$ and $175\,\mathrm{mV}$, an output current of
$120\,\mu\mathrm{A}$, and an output resistance of at least
$25\,\mathrm{k}\Omega$. Determine the width and length of each transistor.

\textbf{Solution.}

\textbf{(a) Process parameters and channel length.}
The single-finger $W/L=5/0.27\,\mu\mathrm{m}$ NFET in the Lab 1
report provides the following extracted parameters:
\[
\mu_nC_{ox}=215.6\,\mu\mathrm{A}/\mathrm{V}^2,\quad
V_{TH}=0.433\,\mathrm{V},\quad \lambda=0.1624\,\mathrm{V}^{-1}.
\]
For $L_1=L_2=0.27\,\mu\mathrm{m}$, the first-order output-resistance
estimate is
\[
R_{\mathrm{out}}\approx\frac{1}{\lambda I_D}
=\frac{1}{0.1624(120\times10^{-6})}
=51.3\,\mathrm{k}\Omega>25\,\mathrm{k}\Omega.
\]

\textbf{(b) Overdrive and transistor widths.}
The center of the specified compliance range gives
\[
V_{OV}\approx V_{\mathrm{compliance}}=0.150\,\mathrm{V},\qquad
V_{GS}\approx V_{TH}+V_{OV}=0.583\,\mathrm{V}.
\]
Neglecting channel-length modulation in the DC current, the
square-law model gives
\[
I_D\approx\frac12\mu_nC_{ox}\frac WL V_{OV}^2,
\qquad
\frac WL=\frac{2(120)}{215.6(0.150)^2}=49.47.
\]
Consequently,
\[
\boxed{W_1=W_2=13.36\,\mu\mathrm{m},\qquad
L_1=L_2=0.27\,\mu\mathrm{m}.}
\]
With $I_{\mathrm{in}}=120\,\mu\mathrm{A}$, equal dimensions give a
nominal 1:1 current ratio. The stated widths are total effective
widths, not individual finger widths.

\textbf{(c) Operating region and interpretation.}
M1 is diode-connected, with $V_{DS1}=V_{GS1}\approx0.583\,\mathrm{V}
>V_{OV1}$. M2 is in saturation for approximately
$V_{\mathrm{out}}\ge0.150\,\mathrm{V}$ under this model.

The dimensions follow from a first-order model whose extracted
parameters depend on geometry and bias. Mobility degradation,
characterized by $\theta=0.9999\,\mathrm{V}^{-1}$ in Lab 1, is omitted
from the sizing equation. Channel-length modulation introduces an
output-voltage dependence into the drain current. The approximate resistance
$1/(\lambda I_D)$ also neglects the finite $(1+\lambda V_{DS})$
correction. These effects require simulation to determine the actual
compliance voltage, output current, and output resistance.

\newpage
\item \textbf{Low-voltage cascode current-mirror design.}
Design a low-voltage cascode current mirror with a 1:2 input-to-output
current ratio, a $60\,\mu\mathrm{A}$ input current, a low-frequency
output resistance greater than $1\,\mathrm{M}\Omega$, and a compliance
voltage below $500\,\mathrm{mV}$. Determine all transistor widths and
lengths and the bias current. Validate saturation at the DC operating point.

\textbf{Solution.}

\textbf{(a) Topology and design assumptions.}
In the topology of Fig.~5-3 in the Lab 5 manual, M1 and M2 are the
lower mirroring devices; M3 and M4 are their respective upper cascode
devices. The drain of M3 connects to the gates of M1 and M2.
The diode-connected bias device $M_B$ supplies the gate voltage of
M3 and M4. The sources of M1, M2, and $M_B$ connect to $V_{SS}$;
all NFET bulk terminals connect to $V_{SS}$.

The design uses the Lab 1 parameters of Question 2, with
\[
L_1=L_2=L_3=L_4=L_B=0.27\,\mu\mathrm{m},\qquad
V_{OV1}=V_{OV2}=V_{OV3}=V_{OV4}=0.150\,\mathrm{V}.
\]
The analysis neglects body effect and mobility degradation. The
square law determines the device dimensions, and $\lambda$ determines
the small-signal output-resistance estimate.

\textbf{(b) Reference and output branch dimensions.}
For M1 and M3, $I_D=60\,\mu\mathrm{A}$:
\[
\left(\frac WL\right)_{1,3}
=\frac{2(60)}{215.6(0.150)^2}=24.74,
\qquad W_1=W_3=6.68\,\mu\mathrm{m}.
\]
For M2 and M4, $I_D=120\,\mu\mathrm{A}$:
\[
\left(\frac WL\right)_{2,4}=49.47,
\qquad W_2=W_4=13.36\,\mu\mathrm{m}.
\]
The doubled output-branch widths provide the nominal 1:2 current ratio.

\textbf{(c) Bias branch and compliance voltage.}
Setting $V_{DS1}\approx V_{DS2}\approx0.250\,\mathrm{V}$ provides
$0.100\,\mathrm{V}$ of nominal saturation margin for the lower devices.
The required cascode gate voltage is
\[
V_B=V_{S3}+V_{GS3}=0.250+(0.433+0.150)=0.833\,\mathrm{V}.
\]
Thus $V_{OVB}=V_B-V_{TH}=0.400\,\mathrm{V}$. For
$I_{\mathrm{bias}}=60\,\mu\mathrm{A}$, the bias-device dimensions are
\[
\left(\frac WL\right)_B
=\frac{2(60)}{215.6(0.400)^2}=3.479,
\qquad \boxed{W_B\approx0.94\,\mu\mathrm{m}.}
\]
The predicted compliance voltage is
\[
\boxed{V_{\mathrm{compliance}}\approx V_{S4}+V_{OV4}
=0.250+0.150=0.400\,\mathrm{V}<0.500\,\mathrm{V}.}
\]

\newpage
\textbf{(d) Saturation verification at the DC operating point.}
At the selected operating point, $V_{\mathrm{out}}=1.000\,\mathrm{V}$.
Ideal current sources supply the input and bias
currents. Since the drain of M3 is connected to the gate of M1,
$V_{D3}=V_{G1}\approx0.583\,\mathrm{V}$. The first-order saturation
condition is $V_{DS}\ge V_{GS}-V_{TH}=V_{OV}$.

\begin{center}
\begin{tabular}{@{}lccc@{}}
\toprule
Device & $V_{DS}$ (V) & $V_{OV}$ (V) & Saturation check\\
\midrule
M1 & 0.250 & 0.150 & $0.250>0.150$\\
M2 & 0.250 & 0.150 & $0.250>0.150$\\
M3 & $0.583-0.250=0.333$ & 0.150 & $0.333>0.150$\\
M4 & $1.000-0.250=0.750$ & 0.150 & $0.750>0.150$\\
$M_B$ & 0.833 & 0.400 & $0.833>0.400$\\
\bottomrule
\end{tabular}
\end{center}
All five devices satisfy the saturation condition under the stated
first-order assumptions.

\textbf{(e) Low-frequency output resistance.}
For the output devices M2 and M4,
\[
r_{o2}\approx r_{o4}\approx\frac{1}{0.1624(120\times10^{-6})}
=51.3\,\mathrm{k}\Omega,
\qquad g_{m4}\approx\frac{2I_D}{V_{OV4}}=1.60\,\mathrm{mS}.
\]
With the cascode gate held at its small-signal bias voltage and
neglecting body effect,
\begin{align*}
R_{\mathrm{out}}&\approx r_{o2}+r_{o4}+g_{m4}r_{o2}r_{o4}\\
&\approx\boxed{4.32\,\mathrm{M}\Omega>1\,\mathrm{M}\Omega}.
\end{align*}

\textbf{(f) Design summary and model validity.}
\begin{center}
\begin{tabular}{@{}lccc@{}}
\toprule
Device & Total $W$ ($\mu$m) & $L$ ($\mu$m) & Nominal $I_D$ ($\mu$A)\\
\midrule
M1 & 6.68 & 0.27 & 60\\
M2 & 13.36 & 0.27 & 120\\
M3 & 6.68 & 0.27 & 60\\
M4 & 13.36 & 0.27 & 120\\
$M_B$ & 0.94 & 0.27 & 60\\
\bottomrule
\end{tabular}
\end{center}
The nominal bias current is $I_{\mathrm{bias}}=60\,\mu\mathrm{A}$.
Because M3 and M4 have elevated source potentials with their bulks
at $V_{SS}$, body effect raises their thresholds and can reduce the
lower devices' saturation margin relative to this calculation.
The Lab 1 summary does not provide body-effect parameters, so its
effect on the operating point remains unquantified. Simulation of
the device operating points determines whether $W_B$,
$I_{\mathrm{bias}}$, or other dimensions require adjustment. The
tabulated values therefore represent analytical design estimates.

\end{enumerate}
\end{document}
